% !TeX root = Svedendhal.tex
\documentclass{article}
\usepackage{amsmath, amssymb, physics, graphicx, bm}

\begin{document}
	
	\title{Derivation of the Modified Fresnel Coefficients}
	\author{Based on Bedeaux and Vlieger's Formalism}
	\date{\today}
	\maketitle
	
	\section{Introduction}
	This section describes the derivation of modified Fresnel coefficients, following the approach of Bedeaux and Vlieger \cite{bedeaux2001optical}. We begin with Ampere's circuital law:
	\begin{equation}
		\nabla \times \vb{H} = \frac{1}{c} \vb{I} + \frac{1}{c} \frac{\partial \vb{D}}{\partial t},
	\end{equation}
	where $\vb{H}$ and $\vb{D}$ are the magnetic and electric displacement fields, $\vb{I}$ is the current density, and $c$ is the speed of light.
	
	\section{Boundary Conditions and Field Decomposition}
	We consider a boundary in the $xy$-plane and decompose the fields:
	\begin{align}
		\vb{H} &= \vb{H}^{+} \theta(z) + \vb{H}^{-} \theta(-z) + \vb{H}^{s} \delta(z), \\
		\vb{D} &= \vb{D}^{+} \theta(z) + \vb{D}^{-} \theta(-z) + \vb{D}^{s} \delta(z),
	\end{align}
	where $\theta(z)$ and $\delta(z)$ are the Heaviside and Dirac delta functions, respectively. The fields $\vb{H}^s$ and $\vb{D}^s$ satisfy $D_{xs} = D_{ys} = H_{xs} = H_{ys} = 0$.
	
	Applying Ampere’s law at $z=0$, we find:
	\begin{equation}
		\nabla \times \vb{H}_s^z + \hat{z} \times (\vb{H}_{||}^{+} - \vb{H}_{||}^{-}) \delta(z) = \frac{1}{c} \vb{I}_s + \frac{1}{c} \frac{\partial \vb{D}_s}{\partial t}.
	\end{equation}
	Multiplying by $(\hat{z} \times)$ and rearranging, we obtain:
	\begin{equation}
		\vb{H}_{||}^{+} - \vb{H}_{||}^{-} = -\hat{z} \times \frac{1}{c} \vb{I}_{||}^s - \hat{z} \times \frac{1}{c} \frac{\partial \vb{P}_{||}^s}{\partial t} - \nabla_{||} M_{z}^s,
	\end{equation}
	where $\vb{P}_s$ and $\vb{M}_s$ represent the electric and magnetic polarization densities.
	
	Defining a generalized displacement field $\vb{D}' = \vb{D} + i/\omega \vb{I}$ and assuming no magnetic polarization at the boundary, we obtain:
	\begin{equation}
		\vb{H}_{||}^{+} - \vb{H}_{||}^{-} = \frac{i \omega}{c} \hat{z} \times \vb{P}_{||}^s.
	\end{equation}
	
	\section{Modified Reflection and Transmission Coefficients}
	For s-polarized light, we assume polarization along $y$, leading to:
	\begin{equation}
		\vb{P}_{||} = -\rho \alpha(\omega) E_{exc, x},
	\end{equation}
	where $\alpha(\omega)$ is the polarizability of a single nanoparticle and $\rho$ is the surface density of nanoparticles. The exciting field is:
	\begin{equation}
		E_{exc, y} = E_t = (E_i + E_r).
	\end{equation}
	By applying boundary conditions for the magnetic field:
	\begin{align}
		H_{x}^{-} &= n_i (-E_i + E_r) \cos \theta_i, \\
		H_{x}^{+} &= -n_t E_t \cos \theta_t = -n_t (E_i + E_r) \cos \theta_t.
	\end{align}
	Inserting these into the boundary condition, we derive the modified reflection coefficient:
	\begin{equation}
		r_s = \frac{n_i \cos \theta_i - n_t \cos \theta_t + i \frac{\omega}{c} \rho \alpha(\omega)}{n_i \cos \theta_i + n_t \cos \theta_t - i \frac{\omega}{c} \rho \alpha(\omega)}.
	\end{equation}
	For transmission, using $t_s = r_s + 1$:
	\begin{equation}
		t_s = \frac{2 n_i \cos \theta_i}{n_i \cos \theta_i + n_t \cos \theta_t - i \frac{\omega}{c} \rho \alpha(\omega)}.
	\end{equation}
	
	The same procedure can be applied for p-polarized light, yielding:
	\begin{equation}
		r_p = \frac{n_t \cos \theta_i - n_i \cos \theta_t - i \frac{\omega}{c} \rho \alpha(\omega) \cos \theta_i \cos \theta_t}{n_t \cos \theta_i + n_i \cos \theta_t - i \frac{\omega}{c} \rho \alpha(\omega) \cos \theta_i \cos \theta_t}.
	\end{equation}
	And for transmission:
	\begin{equation}
		t_p = \frac{2 n_i \cos \theta_i}{n_t \cos \theta_i + n_i \cos \theta_t - i \frac{\omega}{c} \rho \alpha(\omega) \cos \theta_i \cos \theta_t}.
	\end{equation}
	
	\section{Conclusion}
	We have derived the modified Fresnel coefficients by incorporating nanoparticle polarizability and surface density into boundary conditions. This formulation extends classical Fresnel equations and provides a framework for understanding light interaction with structured interfaces.
	
	\bibliographystyle{unsrt}
	\begin{thebibliography}{9}
		\bibitem{bedeaux2001optical} Bedeaux, D., and Vlieger, J. \textit{Optical Properties of Surfaces}. Imperial College Press, 2001.
	\end{thebibliography}
	
	\section{Supporting Information}
	
	\subsection{Derivation of the Modified Fresnel Coefficients}
	This section aims to describe the origin of Eq. (3,4) in the main paper. We follow Bedeaux and Vlieger \cite{bedeaux2001optical}, to which we refer the reader for a more extensive derivation.
	
	To derive the modified Fresnel coefficients, it is convenient to start from Ampere's circuital law, which may be written as:  
	\begin{equation}
		\nabla \times \vb{H} = \frac{1}{c} \vb{I} + \frac{1}{c} \frac{\partial \vb{D}}{\partial t},
	\end{equation}
	where $\vb{H}$ and $\vb{D}$ are the magnetic and electric displacement fields, $\vb{I}$ is the current density, and $c$ is the speed of light. 
	
	Considering a boundary in the $xy$-plane, we may divide the $\vb{H}$ and $\vb{D}$ fields into three components:
	\begin{align}
		\vb{H} &= \vb{H}^{+} \theta(z) + \vb{H}^{-} \theta(-z) + \vb{H}^{s} \delta(z), \\
		\vb{D} &= \vb{D}^{+} \theta(z) + \vb{D}^{-} \theta(-z) + \vb{D}^{s} \delta(z),
	\end{align}
	where $+$, $-$, and $s$ denote the fields above, below, and within the boundary. $\theta(z)$ and $\delta(z)$ are the Heaviside and Dirac delta functions. It is convenient to choose $\vb{D}^{s}$ and $\vb{H}^{s}$ such that $D_{xs} = D_{ys} = H_{xs} = H_{ys} = 0$, without any loss in generality.
	
	Using Ampere’s circuital law at $z = 0$, and noting that $\partial \theta(\pm z) / \partial z = \pm \delta(z)$, it can be shown that:
	\begin{equation}
		\nabla \times \vb{H}^{z}_s + \hat{z} \times (\vb{H}_{||}^{+} - \vb{H}_{||}^{-}) \delta(z) = \frac{1}{c} \vb{I}_s + \frac{1}{c} \frac{\partial \vb{D}_s}{\partial t}.
	\end{equation}
	Multiplying by $(\hat{z} \times)$ and making some simple rearrangements, one finds at $z=0$:
	\begin{equation}
		\vb{H}_{||}^{+} - \vb{H}_{||}^{-} = -\hat{z} \times \frac{1}{c} \vb{I}_{||}^s - \hat{z} \times \frac{1}{c} \frac{\partial \vb{P}_{||}^s}{\partial t} - \nabla_{||} M_{z}^s.
	\end{equation}
	
	Here, $\vb{P}_{||}^{s} = \vb{D}_{||}^{s}$ and $M_{z}^{s} = -H^{z}_{s}$ are the electric and magnetic polarizations, and $\nabla_{||} = \left(\frac{\partial}{\partial x}, \frac{\partial}{\partial y} \right)$. 
	
	By defining a generalized electric displacement field $\vb{D}' = \vb{D} + i/\omega \vb{I}$, the electric current is absorbed in $\vb{D}'$. Assuming no magnetic polarization in the boundary, this leads us to the boundary condition:
	\begin{equation}
		\vb{H}_{||}^{+} - \vb{H}_{||}^{-} = \frac{i \omega}{c} \hat{z} \times \vb{P}_{||}^{s}.
	\end{equation}
	
	For s-polarization, where the polarization is along $y$, we assume:
	\begin{equation}
		\vb{P}_{||} = -\rho \alpha(\omega) E_{exc, x},
	\end{equation}
	where $\alpha(\omega)$ is the polarizability of a single nanoparticle and $\rho$ is the surface density of nanoparticles. The exciting field is:
	\begin{equation}
		E_{exc, y} = E_t = (E_i + E_r).
	\end{equation}
	Applying boundary conditions for the magnetic field:
	\begin{align}
		H_{x}^{-} &= n_i (-E_i + E_r) \cos \theta_i, \\
		H_{x}^{+} &= -n_t E_t \cos \theta_t = -n_t (E_i + E_r) \cos \theta_t.
	\end{align}
	By combining these relationships, we find:
	\begin{equation}
		n_i \cos \theta_i (E_i - E_r) - n_t \cos \theta_t (E_i + E_r) = -\frac{i \omega}{c} \rho \alpha(\omega) (E_i + E_r).
	\end{equation}
	Rearranging, we obtain the modified reflection coefficient:
	\begin{equation}
		r_s = \frac{n_i \cos \theta_i - n_t \cos \theta_t + i \frac{\omega}{c} \rho \alpha(\omega)}{n_i \cos \theta_i + n_t \cos \theta_t - i \frac{\omega}{c} \rho \alpha(\omega)}.
	\end{equation}
	For transmission:
	\begin{equation}
		t_s = \frac{2 n_i \cos \theta_i}{n_i \cos \theta_i + n_t \cos \theta_t - i \frac{\omega}{c} \rho \alpha(\omega)}.
	\end{equation}
	The same procedure applies for p-polarized light, yielding:
	\begin{equation}
		r_p = \frac{n_t \cos \theta_i - n_i \cos \theta_t - i \frac{\omega}{c} \rho \alpha(\omega) \cos \theta_i \cos \theta_t}{n_t \cos \theta_i + n_i \cos \theta_t - i \frac{\omega}{c} \rho \alpha(\omega) \cos \theta_i \cos \theta_t}.
	\end{equation}
	And for transmission:
	\begin{equation}
		t_p = \frac{2 n_i \cos \theta_i}{n_t \cos \theta_i + n_i \cos \theta_t - i \frac{\omega}{c} \rho \alpha(\omega) \cos \theta_i \cos \theta_t}.
	\end{equation}
	
\end{document}

